\section{Protocol details}\label{app:protocol-details}

Algorithm~\ref{alg:policy} states the execution order, including trace eligibility, manifest control, fallback, and both compilation outcomes.
\begin{algorithm}[htbp]
\caption{PACE with extra-cost margin $\epsilon$.
Family identities, cost profiles, routing charges, and the scenario hazard are supplied.}
\label{alg:policy}
\begin{algorithmic}[1]
\State Initialize cumulative charges $K\gets0$ and agent reference $A\gets0$
\For{each arrival $t$}
\State Load the arriving family's costs and state
\State $A\gets A+\tau_0+c$
\State Expire idle entries and tentatively relist the family's stored live program
\State Compute $U_t^{\mathrm{serv}}$ from Equation~\ref{eq:service-reserve}
\If{$K+U_t^{\mathrm{serv}}>(1+\epsilon)A$}
\State Clear the manifest and select agent service
\EndIf
\State Pay routing and serve the task; on a detected program failure, pay one agent fallback
\State Add the routing and serving charges to $K$; log completed agent traces
\State Invalidate the program after a detected break and update the recurrence records
\If{at least $k$ agent runs are complete, no live program exists, and Equation~\ref{eq:trigger} holds}
\If{$K+\max(C,C^{\mathrm{fail}})\leq(1+\epsilon)A$}
\State Compile and verify a candidate; increment $j$
\If{verification passes}
\State $K\gets K+C$; increment $g$; store and list the program
\Else
\State $K\gets K+C^{\mathrm{fail}}$
\EndIf
\EndIf
\EndIf
\EndFor
\end{algorithmic}
\end{algorithm}

\subsection{Compilation spend and program lifetime}
\label{app:reference}
\label{sec:breakeven}
For independent attempts with fixed verification probability $p>0$, the number of failures before a verified program has mean $(1-p)/p$.
The expected compilation spend is therefore
\begin{equation}
B(p)=C+\left(\frac1p-1\right)C^{\mathrm{fail}}.
\label{eq:breakeven-law}
\end{equation}
Substituting the observed-history estimate $\hat p$ gives PACE's proposal price $\hat B$.
This calculation excludes agent service between attempts and assumes a fixed retry distribution, which the repeated GUI measurements do not establish.
The cost guarantee instead counts each actual model charge and does not use this independence assumption.

\label{sec:stationary}
A replacement calculation explains the lifetime cap $1/h$.
Suppose a working program can be obtained at deterministic cost $B$ and every use independently breaks it with probability $h>0$.
A replacement cycle has expected length $1/h$, including its terminal break.
It pays $B$, extraction on every use, ordinary fallback with probability $q$ on nonterminal uses, and one terminal fallback.
Dividing expected cycle cost by expected length gives
\begin{equation}
c^{\mathrm{prog}}+(1-h)qc+h(B+c).
\label{eq:replacement}
\end{equation}
This is below agent cost $c$ exactly when $s_h/h>B$.
At $\hat H=1/h$ with no routing charge, the proposal has the same form, using estimated spend $\hat B$.
Shorter recurrence projections further limit the uses available to repay compilation.
This reference calculation compares two specified policies and does not establish optimal replacement under unknown parameters.

\subsection{Cost of a conservative reservation}
\label{app:reservation-example}
A single-family example has $300$ arrivals, agent cost $1$, successful compilation cost $2$, possible failed cost $200$, and no program-use cost or drift.
The compiler succeeds, but the policy does not know that in advance.
With $\epsilon=0.25$, PACE's extra budget never reaches the $200$ reservation, so it uses the agent and pays $300$.
Eager retry pays $3$ for source runs and $2$ for compilation.
The example shows how reserving the failure branch can prevent a profitable compilation while preserving the agent-relative constraint.
